\section{Experiments}

The first experiment we have done consist in fitting a number of binary trees to a synthetic dataset and inspecting how the gap $\expv{[R(\hat f_D) - \hat R(\hat f_D)]}$ changes as we increase model complexity $\vert\cH\vert$. Specifically, the experiment consists in 
\begin{enumerate}
    \item support: we discretize an interval $I\subset \R$ into $N$ numbers (we chose $I=[0,10], N = 10,000$);
    \item sample $x$: we sample $x$ from the support, uniformly according to a chosen distribution (we start with the uniform one, then inspect a Gaussian one and an adversarial one);
    \item sample $y$: we create the outcome variables according to the law \[y = \sgn(f(x)+\sigma\varepsilon), \qquad \varepsilon\sim\mathcal{N}(0,1), \sigma\in\R\]where $\sigma$ controls the amount of noise (we use values of $1,2,3,5$) and $f:I\to\R$ is a convex function chosen to yield roughly half of the points with positive image and half negative (we use $f(x) = (x-5)^2 -6.25$);
    \item dataset generation: then, we create a training dataset with size $n<<N$ and a test dataset with size $N<<n_\text{test}$ indistinguishable from $\infty$ (we choose $n=10,20,50$ and $n_\text{test}=200,000$);
    \item select depths: we select a list of depths for our trees, specifically each integer in a multiple of the range of the overfitting threshold, which is $\log_2(n)$;
    \item fit models: we fit a number of binary trees with different depths to this data and evaluate for each one the error on the training set ($\hat R$) and the error on the test set ($R$);
    \item plot: we plot the two errors against the depth of the trees, highlighting the overfitting depth. 
\end{enumerate}
what we observe is that the training data size and the standard deviation of the noise have inverse (respectively, positive and negative) influences on the performance of the model on training data, hence we balance them accordingly. Also, for a greater number of training samples size we inspect a greater number of depths (e.g. up to $4*$floor(threshold)) to observe the long-term trend. 

